\chapter{本体构建方法论}
\chapterart{ch03}{本体构建的生命周期循环}

\begin{introbox}
\textbf{【导读】}
第 2 章回答了「本体由什么构成」，本章回答下一个问题：从一张白纸出发，
如何一步步把领域知识变成一个可用、可验收、可维护的本体？
本章给出一条完整的工程路线：先用\textbf{能力问题}（CQ）划定范围、写下验收标准；
再沿从轻到重的\textbf{构建流程}（Ontology 101 七步法与 METHONTOLOGY）把概念逐步落成模型；
借 \textbf{NeOn 场景}与\textbf{本体复用}站上已有成果的肩膀；
用 \textbf{OntoClean} 给类层次做一次质量体检；
最后以\textbf{协作与治理}机制保证多人长期共建不失控。
第 9 章实战系统的设计正是从本章的 CQ 清单开始的。
\end{introbox}

\chapterbinding{semantica.chapter\_packages.vol1.ch03}{partial}{blocked}
{CQ 注册表、CQ1 查询、正例、单因反例及
\texttt{OE-V1-CH03-SCN-CQ-ACCEPTANCE-001} 已原生绑定；方法材料作为包内工程规则保留。}
{Ontology 101/METHONTOLOGY 阶段门禁和 OntoClean 元属性检查器尚未实现；
CQ1 场景通过只证明声明夹具符合 exact oracle，不证明整套方法论已自动化。}

\section{为什么需要方法论：三个失败故事}

讲方法论之前，先讲三个失败的故事。它们都是虚构的，但做过知识工程项目的读者
多半会觉得似曾相识——每个故事对应一种典型的死法，也对应本章一件工具的用武之地。

\textbf{故事一：过度建模。}某制造企业的知识工程小组立项构建「车间设备本体」。
组里有人读过不少本体论文献，立志做出「体系完备」的模型：金属材料按冶金学谱系
细分到牌号，设备按机械原理分出十几层，连「润滑脂的皂基类型」都有了自己的类。
三个月精雕细琢之后是隆重的上线评审，业务方代表问的第一个问题却是：
「哪些设备这周有空，能接这批铝合金件？」全场沉默——本体里有完备的金属谱系，
却没有「空闲」「排程」这些概念。这套模型后来再也没人打开过。
过度建模的根源是\textbf{为建模而建模}：建模者按学科逻辑追求完备，
而使用者只需要回答手头那二十个问题。没有范围约束的建模，勤奋本身就是浪费。

\textbf{故事二：没有验收标准。}另一个团队吸取了教训，紧贴业务建模。
项目计划上写着「第三季度完成工艺本体」，但没有人说得清「完成」长什么样。
于是验收会开成了观感评审：投影仪上滚动着类层次截图，专家们点头说「看起来挺全」，
项目顺利结项。三个月后调度系统接入，才发现「工序之间的先后约束」这一关键关系
从未建过——当初没有任何一条标准能暴露这个缺口。没有验收标准的项目
有两种结局：要么永远「快好了」，要么「验收了却不可用」。
知识工程需要和软件工程一样的东西：\textbf{可执行的验收测试}。

\textbf{故事三：专家缺位。}第三个团队规格齐全、计划严密，只是省掉了一件「麻烦事」：
和车间里的工艺专家反复开会。工程师们以数据库表结构为蓝本闭门造车——
表名变成类名，外键变成对象属性，两周就「完工」了。上线评审时，
干了三十年的工艺师傅只看了一眼就摇头：「我们从来不这样分类设备。」
这套本体在逻辑上无懈可击，却得不到业务方承认。本体的全部价值在于「共享」——
它必须是各方共同承认的概念契约；专家缺位的本体，建得再对也只是工程师的私人作品。

三个故事翻车的环节各不相同——范围、验收、参与——但病根相同：
把本体构建当成凭感觉与才气行事的艺术创作，而不是一项有流程、有产物、有检查点的
工程活动。软件行业在它的「作坊时代」交过同样的学费，然后发展出需求分析、
设计评审、测试、版本控制这些护栏。\textbf{方法论的本质，就是给知识工程装上
软件工程式的流程护栏}：每个阶段有明确的输入与产物，每个产物有可检查的标准，
每次变更有可追溯的记录。护栏不保证你建出伟大的本体，
但能保证你不掉进上面三个坑。

当然，知识工程并不是软件工程的简单换皮，两者有一处深刻差异值得先说破：
软件的需求多半可以由用户单方面给出，而本体的「需求」——领域知识本身——
分散在专家的头脑、制度文件与数据系统之间，连专家自己也未必能系统地说出来。
因此本体方法论比软件方法论更强调两件事：\textbf{知识获取}
（把隐性知识「逼」出来的访谈技艺与文档考古）和\textbf{共同语言}
（让不懂逻辑的专家能够逐行评审的中间产物）。后文会看到，
能力问题与概念化表格正是为这两件事而生的发明。
另一句丑话也说在前头：方法论不是教条，本章给出的所有流程都允许裁剪——
但裁剪应当是想清楚之后的决定，并且写下理由，
而不是「来不及了先跳过」的借口。

本章的工具箱正是按「堵坑」组织的。\textbf{能力问题}（3.2 节）同时堵住范围坑
与验收坑；\textbf{构建流程}（3.3 至 3.5 节）规定步骤与产物，并把专家评审
做成流程的固定环节；\textbf{复用工程}（3.6 节）防止重复造轮子；
\textbf{OntoClean}（3.7 节）给最容易出错的类层次做质检；
\textbf{协作与治理}（3.8 节）让这一切在多人、长期的现实中持续运转。
四套主要方法的整体定位如下表：

\begin{center}
\begin{tabularx}{\linewidth}{@{}lllX@{}}
\toprule
\textbf{方法} & \textbf{提出} & \textbf{分量} & \textbf{定位与适用场景} \\
\midrule
Ontology 101 & 2001 & 轻 & 七步迭代流程；教学、中小型本体、快速原型 \\
METHONTOLOGY & 1997 & 重 & 完整生命周期与标准产物；企业级交付、需文档与评审 \\
NeOn & 2009 & 中 & 九种可组合场景；从已有标准、词表、本体出发的构建 \\
OntoClean & 2002 & 工具 & 不是流程，而是类层次的元属性质检方法 \\
\bottomrule
\end{tabularx}
\end{center}

这四件工具不互斥，而是互补：典型的项目用 CQ 开局，按七步法或 METHONTOLOGY 推进，
中途按 NeOn 场景引入复用资源，节点处用 OntoClean 体检，全程接受治理机制约束。
下面逐一展开。

给不同处境的读者三条阅读路径。正要启动第一个本体项目的读者，
请按 3.2、3.3、3.7、3.8 的顺序精读，3.4 与 3.5 浏览即可；
接手半途项目的读者，先用 3.2 节给它补一份 CQ 清单、
再用 3.7 节给现有类层次做个体检，诊断清楚再决定续建还是重构；
任务是把企业已有分类标准「本体化」的读者，3.5 与 3.6 两节就是为你写的。
无论走哪条路径，3.9 节的思考题都值得动笔——方法论这种东西，
不亲手用过一遍，永远只是别人的经验。

\section{先问问题，再建模型：能力问题}

回到故事一的评审会现场。如果三个月前，团队先把业务方那天问出的问题
一条条写在纸上——「哪些设备这周有空」「哪些设备能加工铝合金」——
再围绕这张问题清单建模，结局会完全不同：金属谱系不会细到皂基类型，
「空闲」「排程」也绝不会漏掉。这张清单，就是本体工程中最重要的需求工具：
\textbf{能力问题}（Competency Question，CQ）。

\subsection{定义与三重作用}

能力问题是\textbf{本体建成后必须能够通过查询或推理回答的问题}。
这个概念由 Gr\"uninger 与 Fox 在 1990 年代中期的 TOVE 企业本体项目中提出，
此后成为几乎所有本体方法论的共同起点。一句话概括它的双重身份：
CQ 既是本体的\textbf{需求规格说明书}，又是它的\textbf{验收测试用例}。
展开来说，CQ 起三重作用。

\textbf{第一重：范围边界。}CQ 覆盖不到的概念，不要建。
这条纪律直接堵住故事一的坑：如果没有任何 CQ 问到财务折旧，
本体里就不应该出现折旧概念——无论它在领域教科书里多么重要。
范围蔓延在知识工程中比在软件工程中更隐蔽，因为「把概念体系补完整」
听起来总是正当的；CQ 给了团队一个可以理直气壮说「不建」的依据。

\textbf{第二重：设计驱动。}每个 CQ 都能机械地反推出需要的类、属性与公理。
「哪些设备可以加工铝合金」一句话里就藏着两个类（\texttt{Equipment}、
\texttt{Material} 及其子类铝合金）、一条对象属性（\texttt{canProcess}）
和一条子类公理。把全部 CQ 反推完毕，第一版概念清单就自动出现了——
建模不再从灵感开始，而是从需求开始。3.2.5 小节专门演示这套反推。

\textbf{第三重：验收标准。}每条 CQ 都可以翻译成一条对本体的查询
（通常是第 4 章将要正式登场的 SPARQL），跑出的结果与领域专家事先给出的
标准答案一致，即告通过。「建完」从此有了客观定义：\textbf{全部 CQ 验收通过}。
这直接堵住故事二的坑。

实践中还有一个常被问到的细节：CQ 要写到多具体？
「哪些设备可以加工铝合金」与「哪些设备可以加工钛合金」是一条 CQ 还是两条？
惯例是把它们看作同一个\textbf{问题模板}的不同实例化——模板里的「铝合金」
是参数位。写清单时登记模板并配一两个典型参数即可，验收时再用具体参数执行。
这样既避免清单被同型问题灌水，也提醒设计者：本体要支撑的是「这一类问题」，
而不是某个写死的问句。

\subsection{好的能力问题长什么样}

不是随手写下的问题都配叫 CQ。实践中好的 CQ 有三个特征，可以当作入池前的体检单。

\textbf{特征一：具体可查。}「了解车间设备情况」不是 CQ，
「当前哪些设备处于空闲状态」才是。判别方法很简单：闭上眼睛想象
它变成一条查询语句、返回一张结果表——想象不出来，就还不够具体。
模糊的愿望要拆解成若干条可查询的问题，拆不动的部分往往根本不是本体的职责。

\textbf{特征二：有判据。}每条 CQ 必须存在可事先给出的标准答案或判定规则，
否则将来无从验收。「3 号车床现在能不能接这个急单」是好问题，
前提是专家能说清判定逻辑（空闲、能力匹配、无冲突）；
「哪台设备最好用」就缺乏判据——除非先把「好用」定义成可计算的指标。

\textbf{特征三：业务方背书。}每条 CQ 应当来自真实业务场景，
并且有具体的角色（调度员、工艺员、设备科长）愿意确认「这就是我要问的问题」、
愿意为标准答案签字。背书把故事三的坑提前堵上：从需求一开始，
专家就被拉进了流程，而不是等模型建完再来「围观」。

数量上，包内教学材料给出的经验是\textbf{每轮迭代 10～30 条为宜}：
太少撑不起一个像样的范围，太多则说明范围失控或粒度过细，应当分轮消化。

这三个特征也解释了 CQ 与传统需求文档的差别。需求文档说
「系统应支持设备管理」，CQ 说「当前哪些设备处于空闲状态」——
前者是愿景，后者是测试。把愿景翻译成测试的过程会暴露大量含混之处：
「支持」是什么意思？「设备管理」要管到哪一层？这些追问在写 CQ 时
就被迫回答，而不是留到验收现场吵架。所以有经验的团队干脆把 CQ 写作
当成需求澄清工具来用：能顺利写成 CQ 的需求，才是真需求。

\subsection{从专家访谈中提炼能力问题}

CQ 的主要来源是领域专家访谈。但要注意：\textbf{专家不会直接说出 CQ}——
他们说出的是抱怨、流程和故事，提炼是工程师的功课。访谈时与其问
「您需要本体回答什么问题」（多数专家无从答起），不如使用一组久经考验的提问模板：

\begin{itemize}
  \item 「您每天上班第一件事要查什么？」——挖出高频的基础查询。
  \item 「排产（或诊断故障）时，您先看哪些数据，再做什么判断？」——挖出关联推理链。
  \item 「哪类错误一旦发生就很麻烦，最好能提前拦住？」——挖出约束校验需求。
  \item 「月底向上汇报时，您要算哪些数字？」——挖出统计分析需求。
  \item 「新人最常追着您问什么？」——挖出还没沉淀成文档的隐性知识。
\end{itemize}

提炼遵循一个三步小循环：\textbf{记录原话}（保留专家的业务用语，不要当场翻译成
建模术语）、\textbf{改写成可查询形式}（套上「哪些……」「……是什么」的问句外壳，
补全隐含条件）、\textbf{回读确认}（把改写后的问题念给专家听，请他确认
「问的还是原来那件事」）。回读这一步同时完成了业务方背书，一举两得。
除了访谈，报表、看板、巡检单这些纸面制品也是 CQ 富矿——
每一个表头背后，都站着一条已经被业务日常验证过的潜在 CQ。

访谈还有三条纪律。其一，\textbf{追实例，不追定义}：别问「什么是瓶颈设备」，
要问「上个月哪台设备成了瓶颈、您当时是怎么发现的」——专家给出的具体判断过程
远比抽象定义可靠，定义可以事后归纳。其二，\textbf{记录分歧，不当场调解}：
两位专家对同一术语的用法不一致时，原样记下两种用法——这往往预示着将来
要建两个概念或一对易混淆的兄弟类；当场和稀泥只会把分歧埋进模型，
等它在推理结果里再次现身时就难查多了。其三，\textbf{控制时长与频次}：
专家的耐心是消耗品，一次一小时、多次小步，远比一次榨干三小时可持续。

\subsection{能力问题分类学}

收集来的 CQ 不是一盘散沙。按「回答它需要动用什么机制」分类，
绝大多数 CQ 落进四个篮子：

\begin{center}
\begin{tabularx}{\linewidth}{@{}lXXl@{}}
\toprule
\textbf{类别} & \textbf{典型问句} & \textbf{动用的机制} & \textbf{验收方式} \\
\midrule
基础查询 & 工作单元 A 中有哪些设备？ & 类层次与属性断言的直接匹配 & 对照台账 \\
关联推理 & 任务 T001 与哪些任务冲突？ & 属性链与规则推理（第 5 章） & 专家标准答案 \\
约束校验 & 哪些高功率设备未配冷却？ & 公理约束与反例检测 & 期望空集 \\
统计分析 & 各单元设备利用率是多少？ & 查询聚合（第 4 章） & 对照报表口径 \\
\bottomrule
\end{tabularx}
\end{center}

分类的用处有三：一是\textbf{排难度梯度}——基础查询类最容易实现，适合第一轮迭代
先拿下，建立信心；二是\textbf{定技术选型}——关联推理类提示将来需要规则
（第 5 章 SWRL），统计分析类提示需要查询层聚合而非本体公理；
三是\textbf{切迭代}——每轮迭代认领一小撮 CQ，做完即验收。
下面是贯穿本书的制造领域 CQ 清单，来自 ch03 包内教学资产：

\begin{codebox}
\codeline{\textcolor{cmtgray}{\#\ ==============================}}
\codeline{\textcolor{cmtgray}{\#\ 2.\ 制造领域CQ清单（节选）}}
\codeline{\textcolor{cmtgray}{\#\ ==============================}}
\codeline{\textcolor{cmtgray}{\#\ 基础查询类：}}
\codeline{\textcolor{cmtgray}{\#\ \ \ CQ1\ \ 哪些设备可以加工铝合金？}}
\codeline{\textcolor{cmtgray}{\#\ \ \ CQ2\ \ 当前哪些设备处于空闲状态？}}
\codeline{\textcolor{cmtgray}{\#\ \ \ CQ3\ \ 工作单元A中有哪些设备？}}
\codeblank{}
\codeline{\textcolor{cmtgray}{\#\ 关联推理类：}}
\codeline{\textcolor{cmtgray}{\#\ \ \ CQ4\ \ 任务T001与哪些已排产任务存在时间冲突？}}
\codeline{\textcolor{cmtgray}{\#\ \ \ CQ5\ \ 产品P001应推荐什么加工工艺？}}
\codeline{\textcolor{cmtgray}{\#\ \ \ CQ6\ \ 工单W001的工序先后顺序是什么？}}
\codeblank{}
\codeline{\textcolor{cmtgray}{\#\ 约束校验类：}}
\codeline{\textcolor{cmtgray}{\#\ \ \ CQ7\ \ 哪些高功率设备没有配置冷却系统？（违反公理）}}
\codeline{\textcolor{cmtgray}{\#\ \ \ CQ8\ \ 哪些设备的维护记录超期？}}
\codeblank{}
\codeline{\textcolor{cmtgray}{\#\ 统计分析类：}}
\codeline{\textcolor{cmtgray}{\#\ \ \ CQ9\ \ 各工作单元的设备利用率是多少？}}
\codeline{\textcolor{cmtgray}{\#\ \ \ CQ10\ 过去一个月哪类故障发生频率最高？}}
\codeblank{}
\codeline{\textcolor{cmtgray}{\#\ 注：CQ1/CQ2/CQ4/CQ5/CQ6\ 在第9章实战系统中有完整实现}}
\codeblank{}
\end{codebox}

\assetsource{ch03}{competency-questions}

请记住其中的 CQ1 至 CQ5——第 9 章沿用这些编号，但当前只有 CQ1 绑定了
正例、单因反例与 exact oracle；其余查询即使已迁入，也不能在缺少夹具与 oracle 时假绿。
还要注意约束校验类的特殊气质：CQ7 问的是「哪些设备\textbf{违反了}应有的配置」，
这类问题的理想答案是「一个都没有」——这个细节在 3.2.6 小节的验收设计中将变得关键。

四类问题对本体的「压强」也不相同。基础查询类几乎只消耗类层次与属性断言，
是检验骨架立不立得住的轻负载测试；关联推理类开始牵动属性链与规则，
还会逼出隐藏的中间概念——CQ4 就逼出了「时间区间」与「冲突」这两个
访谈中没人主动提起的类；约束校验类要求把业务红线写成公理，
是本体从「资料库」升级为「裁判员」的标志；统计分析类则提醒我们本体不是孤岛——
聚合计算发生在查询层，本体的职责是把口径定义清楚：什么算「利用」、
哪段时间算「一个月」。一份健康的 CQ 清单应当四类俱全，
否则建出来的本体注定偏科。

\subsection{从能力问题反推本体元素}

CQ 的第二重作用——设计驱动——值得单独演示。反推有一套近乎机械的语法线索：
问句中的\textbf{名词}对应类或实例（「设备」「铝合金」），
\textbf{动词与介词短语}对应属性（「加工」对应 \texttt{canProcess}，
「位于」对应 \texttt{locatedIn}），\textbf{种属修饰}对应子类公理
（「铝合金是一种材料」），\textbf{时间与数值条件}对应数据属性与比较逻辑。
包内 \texttt{competency-questions} 教学资产对三条代表性 CQ 做了完整分解：

\begin{codebox}
\codeline{\textcolor{cmtgray}{\#\ ==============================}}
\codeline{\textcolor{cmtgray}{\#\ 3.\ CQ\ \(\rightarrow\)\ 本体元素的反推}}
\codeline{\textcolor{cmtgray}{\#\ ==============================}}
\codeline{\textcolor{cmtgray}{\#\ 每个CQ可机械地分解出本体需要的元素：}}
\codeblank{}
\codeline{\textcolor{cmtgray}{\#\ CQ1\ "哪些设备可以加工铝合金？"}}
\codeline{\textcolor{cmtgray}{\#\ \ \ 类：\ \ \ \ Equipment（设备）、Material（材料）、AluminumAlloy（铝合金）}}
\codeline{\textcolor{cmtgray}{\#\ \ \ 属性：\ \ canProcess（对象属性：Equipment\ \(\rightarrow\)\ Material）}}
\codeline{\textcolor{cmtgray}{\#\ \ \ 公理：\ \ AluminumAlloy\ \(\sqsubseteq\)\ Material}}
\codeline{\textcolor{cmtgray}{\#\ \ \ 实例：\ \ 具体设备与材料个体}}
\codeblank{}
\codeline{\textcolor{cmtgray}{\#\ CQ4\ "任务T001与哪些任务存在时间冲突？"}}
\codeline{\textcolor{cmtgray}{\#\ \ \ 类：\ \ \ \ Task、TimeInterval、Conflict}}
\codeline{\textcolor{cmtgray}{\#\ \ \ 属性：\ \ hasTimeInterval、assignedTo、conflictsWith}}
\codeline{\textcolor{cmtgray}{\#\ \ \ 规则：\ \ 同设备时间区间重叠\ \(\rightarrow\)\ 冲突（SWRL实现见第5章）}}
\codeblank{}
\codeline{\textcolor{cmtgray}{\#\ CQ8\ "哪些设备的维护记录超期？"}}
\codeline{\textcolor{cmtgray}{\#\ \ \ 类：\ \ \ \ MaintenanceRecord}}
\codeline{\textcolor{cmtgray}{\#\ \ \ 属性：\ \ hasMaintenanceDate（数据属性：xsd:dateTime）、maintenanceCycle}}
\codeline{\textcolor{cmtgray}{\#\ \ \ 约束：\ \ 当前时间\ -\ 上次维护时间\ >\ 周期\ \(\rightarrow\)\ 超期（需注入当前时间）}}
\codeblank{}
\end{codebox}
\color{black}

\assetsource{ch03}{competency-questions}

把全部 CQ 的反推结果汇总、去重，就得到第一版\textbf{类清单、属性清单与公理清单}——
这正是 3.3 节七步法进入「定义类层次」之前最好的输入。反推还有一个副产品：
提前暴露能力边界。CQ8 的分解显示，「维护是否超期」需要拿\textbf{当前时间}
与维护日期做算术比较，而这超出了纯 OWL 推理的能力范围，
必须放到查询层或规则层解决（第 5 章详述）。在纸面上发现这一点，
远好于在联调现场发现。

团队作业时，可以把反推结果整理成一张\textbf{需求—元素矩阵}：
行是 CQ，列是类、属性与公理，命中处打勾。这张矩阵有两个妙用。
纵向看：被多条 CQ 共用的元素就是骨干，值得多花打磨功夫；
横向看：勾迟迟补不齐的那条 CQ，就是当前迭代的风险点。
矩阵还能反向审计——某个类若在所有行都落空，要么它是复用顶层带来的结构件，
要么就是该删的「为建模而建模」遗迹。故事一里的皂基类型，
在这张矩阵上一行勾都得不到。

\subsection{把验收自动化：从 CQ 到 SPARQL 回归}

验收标准要想长出牙齿，必须\textbf{可判定}。每条 CQ 不只配查询，还要绑定适用范围、
正例、单因反例与 exact oracle；本体每次变更后由同一 runner 重跑。先看两个教学例子：

\begin{codebox}
\codeline{\textcolor{cmtgray}{\#\ CQ1\ 的验收查询：}}
\codeline{\textcolor{cmtgray}{\#}}
\codeline{\textcolor{cmtgray}{\#\ PREFIX\ mfg:\ <http://example.org/manufacturing\#>}}
\codeline{\textcolor{cmtgray}{\#\ SELECT\ ?equipment\ WHERE\ \{}}
\codeline{\textcolor{cmtgray}{\#\ \ \ \ \ ?equipment\ rdf:type/rdfs:subClassOf*\ mfg:Equipment\ .}}
\codeline{\textcolor{cmtgray}{\#\ \ \ \ \ ?equipment\ mfg:canProcess\ ?m\ .}}
\codeline{\textcolor{cmtgray}{\#\ \ \ \ \ ?m\ rdf:type\ mfg:AluminumAlloy\ .}}
\codeline{\textcolor{cmtgray}{\#\ \}}}
\codeline{\textcolor{cmtgray}{\#}}
\codeline{\textcolor{cmtgray}{\#\ 验收判据：返回结果与领域专家给出的标准答案一致}}
\codeblank{}
\codeline{\textcolor{cmtgray}{\#\ CQ7\ 的验收查询（约束校验类，期望结果为空集）：}}
\codeline{\textcolor{cmtgray}{\#}}
\codeline{\textcolor{cmtgray}{\#\ SELECT\ ?e\ WHERE\ \{}}
\codeline{\textcolor{cmtgray}{\#\ \ \ \ \ ?e\ rdf:type\ mfg:HighPowerEquipment\ .}}
\codeline{\textcolor{cmtgray}{\#\ \ \ \ \ FILTER\ NOT\ EXISTS\ \{\ ?e\ mfg:requiresCooling\ ?c\ \}}}
\codeline{\textcolor{cmtgray}{\#\ \}}}
\codeline{\textcolor{cmtgray}{\#}}
\codeline{\textcolor{cmtgray}{\#\ 验收判据：结果非空说明数据或公理有缺陷，需修复后复测}}
\end{codebox}

\assetsource{ch03}{competency-questions}

第一条对应 CQ1。值得留意路径表达式 \texttt{rdf:type/rdfs:subClassOf*}：
它把「数控车床的实例」也算进「设备」，否则挂在子类下的设备会漏网——
这是新手最常踩的查询陷阱（SPARQL 语法在第 4 章系统展开，这里先混个脸熟）。
验收判据不是“看起来差不多”，而是变量、行数与多重集和领域专家冻结的标准答案一致。
当前 ch03 package 已为 CQ1 原生绑定这一组证据。

第二条对应 CQ7，设计哲学完全不同：\textbf{约束校验类 CQ 的期望结果可以是空集}。
查询专门去找「高功率却没有配置冷却」的反例，返回任何一行都意味着
数据或公理有缺陷——结果非空本身就是告警。一正一反两类查询，
合起来展示了“正向答案”与“找反例”两种测试形态；但 CQ7 当前没有章内执行场景，
不能仅凭教学查询声称已验收。

最后一步是把测试装进 Semantica package：CQ registry、查询、夹具、场景与 oracle
共同进入 manifest 和内容哈希；每次变化后由 runner 自动回归全部\textbf{声明为必需}
的场景。只有 package release gate 完成才允许发布，不能在书旁另建 CQ 目录或把
“若干查询通过”折算成全绿。
有了这张回归测试网，本体才敢放心重构：改动类层次时，
任何被破坏的旧能力都会在几秒内亮红灯，而不是三个月后在业务现场爆雷。
新需求到来时也不慌——新需求就是新 CQ，入池、反推、建模、验收，进入下一轮循环。

落地这套回归，还有两件小工具值得一提。一是给每条 CQ 一个\textbf{稳定编号}
（CQ1、CQ2……），让评审纪要、变更日志与查询文件可以互相引用——
本书从本章到第 9 章沿用同一套编号，就是这个习惯的示范。
二是给每条验收查询附上\textbf{标准答案的来源与日期}：
答案是哪位专家给的、依据哪天的台账，将来业务变化导致答案过期时，
才知道找谁更新。验收体系自身也需要维护——这件事本身，
就是 3.8 节治理机制的一部分。

\section{七步法精讲：Ontology 101}

手里攥着一沓通过体检的 CQ，终于可以动手建模了。第一套流程选最轻的：
Noy 与 McGuinness 于 2001 年发表的《Ontology Development 101》，
出自斯坦福 Protégé 团队之手，至今仍是流传最广的入门方法论。
它把构建过程拆成七步：确定范围 \(\rightarrow\) 考虑复用 \(\rightarrow\) 枚举术语
\(\rightarrow\) 定义类层次 \(\rightarrow\) 定义属性 \(\rightarrow\) 定义约束
\(\rightarrow\) 创建实例。它的「轻」体现在不强制任何文档产物，
随做随改；它的另一关键词是「循环」——七步从来不是瀑布，这一点留到 3.3.8 小节展开。
本节以贯穿全书的车间设备本体为例逐步精讲，
包内 \texttt{ontology-101-process} 资产给出了完整走查记录。

动手之前先对七步建立整体感：前三步在\textbf{收集}——圈范围、找现货、攒词汇；
中间三步在\textbf{结构化}——搭层次、连关系、上规矩；最后一步\textbf{实例化}。
三段式的重心一路下沉，从领域对话沉到形式模型，再沉到具体数据；
难度的峰值出现在第四步，返工的高发区则在第五、六步之间。
心里有了这张地形图，中途迷路时才知道自己身在何处。

\subsection{第一步：确定领域与范围}

一切从四个问题开始。它们看似朴素，却分别钉死了项目的四个自由度：

\begin{codebox}
\codeline{\textcolor{cmtgray}{\#\ ==============================}}
\codeline{\textcolor{cmtgray}{\#\ Step\ 1.\ 确定领域与范围}}
\codeline{\textcolor{cmtgray}{\#\ ==============================}}
\codeline{\textcolor{cmtgray}{\#\ 用四个问题界定：}}
\codeline{\textcolor{cmtgray}{\#\ \ \ 本体覆盖什么领域？\ \ \ \ \ \ \ \ ——\ 离散制造车间的设备与加工能力}}
\codeline{\textcolor{cmtgray}{\#\ \ \ 本体用来做什么？\ \ \ \ \ \ \ \ \ \ ——\ 设备调度、能力匹配、冲突检测}}
\codeline{\textcolor{cmtgray}{\#\ \ \ 本体要回答哪些问题？\ \ \ \ \ \ ——\ 能力问题清单（见\ competency-questions.txt）}}
\codeline{\textcolor{cmtgray}{\#\ \ \ 谁使用和维护本体？\ \ \ \ \ \ \ \ ——\ 工艺工程师建模，调度系统消费}}
\codeblank{}
\codeline{\textcolor{cmtgray}{\#\ 范围控制原则：CQ覆盖不到的概念不建模}}
\codeline{\textcolor{cmtgray}{\#\ （例：本体不需要覆盖财务折旧，除非有对应CQ）}}
\codeblank{}
\end{codebox}
\color{black}

\assetsource{ch03}{ontology-101-process}

第一问「覆盖什么领域」圈定知识的疆界；第二问「用来做什么」决定建模的侧重——
同样是设备，为调度服务的本体强调能力与状态，为维修服务的本体强调结构与故障模式；
第三问「回答哪些问题」直接挂接 3.2 节的 CQ 清单，是范围的可执行表述；
第四问「谁使用、谁维护」最常被忽略，却决定着形式化深度与工具选择——
供调度系统机器消费的本体必须严格公理化，只供人浏览的词表则可以从轻。
范围控制的纪律仍然是那一条：\textbf{CQ 覆盖不到的概念不建模}。
示例里特意点了财务折旧的名：除非真有一条 CQ 问到它，
否则它再「重要」也进不了这个本体。

还有一个反向的坑要点名：\textbf{让手头数据决定范围}。
「台账里有什么字段就建什么」听起来务实，实则把第四问的答案
让渡给了历史遗留系统——台账没记录的能力关系（哪台设备能加工什么材料）
恰恰可能是 CQ 最需要的知识。正确的顺序永远是：CQ 定范围，
范围定数据需求，缺的数据设法补采；而不是反过来削足适履，
让昨天的数据库决定明天的知识边界。

\subsection{第二步：考虑复用已有本体}

写下第一个类之前，先抬头看看世界——你要建的概念，多半有人建过。
检索顺序遵循「自顶向下」：先看\textbf{顶层本体}，BFO（已成为国际标准
ISO/IEC 21838-2）提供物理对象、过程、质量这些最根本的范畴区分；
再看\textbf{领域本体}，工业本体基础 IOF 给出了工业过程、设备、能力等核心概念；
然后是\textbf{专项本体}，QUDT 管单位与量纲（功率千瓦、长度毫米），
W3C Time Ontology 管时间区间与先后关系——CQ4 的「时间冲突」检测将直接受益。

复用的收益远不止省工。对齐公共顶层，意味着你的本体与同样对齐该顶层的
其他本体天然可互操作；成熟本体里沉淀的建模决策（过程与参与者如何关联、
能力与实现如何区分），替你避开了前人踩过的坑。接入方式此处先记两个名字——
\texttt{owl:imports} 全量引入与 \texttt{rdfs:subClassOf} 概念对齐——
完整的「找、评、用」工程与陷阱清单是 3.6 节的主题，全部资源入口见附录 A。

为什么复用排在第二步，而不是建完再说？因为复用决策影响一切后续步骤：
决定挂接 BFO，意味着你的「过程」「对象」切分要服从它的范畴体系，
第四步的顶层结构因此基本定型；决定引入 QUDT，意味着第五步的数值属性
要带单位建模，而不是裸浮点数。复用是地基选择，地基浇完再换，
代价就是推倒重来。哪怕最终决定完全自建，第二步的检索也绝不白做——
认真读过几个同领域本体再动手的人，对「坑在哪里」的预判完全不同。

\subsection{第三步：枚举重要术语}

现在放下结构包袱，做一件纯体力活：把领域里的重要词汇穷举出来。

\begin{codebox}
\codeline{\textcolor{cmtgray}{\#\ ==============================}}
\codeline{\textcolor{cmtgray}{\#\ Step\ 3.\ 枚举重要术语}}
\codeline{\textcolor{cmtgray}{\#\ ==============================}}
\codeline{\textcolor{cmtgray}{\#\ 不分类、不纠结，先穷举（来自访谈记录与工艺文件）：}}
\codeline{\textcolor{cmtgray}{\#\ \ \ 设备、车床、数控车床、铣床、加工中心、机器人、AGV、}}
\codeline{\textcolor{cmtgray}{\#\ \ \ 材料、铝合金、钛合金、45号钢、}}
\codeline{\textcolor{cmtgray}{\#\ \ \ 工序、车削、铣削、钻孔、热处理、检验、}}
\codeline{\textcolor{cmtgray}{\#\ \ \ 工作单元、产线、车间、}}
\codeline{\textcolor{cmtgray}{\#\ \ \ 功率、转速、精度、利用率、状态、}}
\codeline{\textcolor{cmtgray}{\#\ \ \ 加工、位于、需要、先于、操作……}}
\codeblank{}
\end{codebox}
\color{black}

\assetsource{ch03}{ontology-101-process}

这一步的口诀是\textbf{「不分类、不纠结、先穷举」}——此刻唯一会犯的错误是漏。
不要停下来争论「数控车床该挂在哪」，那是第四步的事；现在只管往清单上添词。
素材来源：专家访谈记录、工艺规程文件、报表表头、信息系统的数据字典，
以及 CQ 文本本身。一个容易忽视的提醒：\textbf{动词也要收}。
「加工」「位于」「先于」「操作」眼下只是清单上的词，
到第五步它们将转生为对象属性。这张平面术语表是后续两步的全部原料，
也是 3.4 节 METHONTOLOGY 术语表的雏形。

如果团队人多，枚举可以开成卡片工作坊：人手一叠卡片，
十分钟内各自写下能想到的术语，然后贴墙归堆。重复出现的词
天然就是骨干候选；只被一个人写出的词则要追问——是个人黑话，
还是真有其事？工作坊的副产品同样宝贵：哪些词当场引发了
「你说的 X 跟我说的不是一回事」，记下来，它们注定是未来评审会的常客。

\subsection{第四步：定义类与类层次}

七步之中，这一步最考验功力——第 2 章讲过，类层次（TBox）是本体的承重结构。
入手有三种策略：

\begin{codebox}
\codeline{\textcolor{cmtgray}{\#\ ==============================}}
\codeline{\textcolor{cmtgray}{\#\ Step\ 4.\ 定义类与类层次}}
\codeline{\textcolor{cmtgray}{\#\ ==============================}}
\codeline{\textcolor{cmtgray}{\#\ 三种策略：}}
\codeline{\textcolor{cmtgray}{\#\ \ \ 自顶向下：Equipment\ \(\rightarrow\)\ ProcessingEquipment\ \(\rightarrow\)\ CNCMachine\ \(\rightarrow\)\ CNCLathe}}
\codeline{\textcolor{cmtgray}{\#\ \ \ 自底向上：先有\ CNCLathe/CNCMill，再归纳出\ CNCMachine}}
\codeline{\textcolor{cmtgray}{\#\ \ \ 混合法（推荐）：先定中层骨干概念，向上挂顶层、向下细化}}
\codeblank{}
\codeline{\textcolor{cmtgray}{\#\ 本例类层次（节选）：}}
\codeline{\textcolor{cmtgray}{\#\ Equipment\ \(\sqsubseteq\)\ PhysicalObject}}
\codeline{\textcolor{cmtgray}{\#\ \ \ ProcessingEquipment\ \(\sqsubseteq\)\ Equipment}}
\codeline{\textcolor{cmtgray}{\#\ \ \ \ \ CNCMachine\ \(\sqsubseteq\)\ ProcessingEquipment}}
\codeline{\textcolor{cmtgray}{\#\ \ \ \ \ \ \ CNCLathe\ \(\sqsubseteq\)\ CNCMachine}}
\codeline{\textcolor{cmtgray}{\#\ \ \ \ \ \ \ CNCMillingMachine\ \(\sqsubseteq\)\ CNCMachine}}
\codeline{\textcolor{cmtgray}{\#\ \ \ MeasurementEquipment\ \(\sqsubseteq\)\ Equipment}}
\codeline{\textcolor{cmtgray}{\#\ Material\ \(\sqsubseteq\)\ PhysicalObject}}
\codeline{\textcolor{cmtgray}{\#\ \ \ MetalMaterial\ \(\sqsubseteq\)\ Material}}
\codeline{\textcolor{cmtgray}{\#\ \ \ \ \ AluminumAlloy\ \(\sqsubseteq\)\ MetalMaterial}}
\codeblank{}
\codeline{\textcolor{cmtgray}{\#\ 常见错误自查：}}
\codeline{\textcolor{cmtgray}{\#\ \ \ 单一子类（某类只有一个子类）\(\rightarrow\)\ 层次可能多余}}
\codeline{\textcolor{cmtgray}{\#\ \ \ 类名用复数\ \(\rightarrow\)\ 统一用单数（CNCLathe\ 而非\ CNCLathes）}}
\codeline{\textcolor{cmtgray}{\#\ \ \ 把实例当类（"3号车床"是个体，不是类）}}
\codeblank{}
\end{codebox}
\color{black}

\assetsource{ch03}{ontology-101-process}

\textbf{自顶向下}从最抽象的概念出发逐层细分。它适合领域里已有公认分类体系
（例如设备分类国家标准）的场景，骨架一步到位；风险是过早抽象——
按学科逻辑造出一堆业务上根本用不到的中间层，正是故事一的病灶之一。
\textbf{自底向上}从台账里的具体概念出发向上归纳，适合数据驱动的场景，
每个类都有实例兜底、绝不悬空；风险是见树不见林，归纳出的层次七零八落。
\textbf{混合法}是多数项目的答案：先抓\textbf{中层骨干概念}——
专家日常口语里最常出现的那些词（车床、铣床、机器人、工序），
它们粒度适中、共识最强、最稳定；然后向上挂接顶层本体，向下按 CQ 需要细化。

层次初成后，用三条自查规则扫一遍，每条都对应一种高发病：

\textbf{规则一：警惕单一子类。}某个类若只有一个子类，这层多半是多余的——
要么补齐它的兄弟类，要么把这一层删掉。这种「独生子」结构往往说明
你在借类层次表达别的东西（版本、别名、部门归属），而那些都另有正道。

\textbf{规则二：类名一律用单数。}\texttt{CNCLathe} 而非 \texttt{CNCLathes}。
类名指称概念本身而非其外延集合；统一单数还能让命名、对齐与代码生成
免于「单复数混战」。

\textbf{规则三：分清实例与类。}「3 号车床」是个体，不是类。
两条实用判据：其一，它还会不会再被细分出成员？不会，就是实例；
其二，你想谈论的是「它本身」（它的功率、它的位置）还是「这一类东西的共性」？
前者是实例，后者才是类。把实例误建成类，是初学者错误榜上的常年冠军。

最后泼一盆冷水：通过了三条自查，类层次也只是「没有明显硬伤」。
更深层的概念混淆——比如把「角色」当成了「类型」——要靠 3.7 节的
OntoClean 体检才能查出来。

另有一针预防剂此处必须先打：\textbf{多重继承要节制}。
OWL 允许一个类有多个父类，但「数控车床既是切削设备、又是数控设备、
又是大型设备」这样的网状层次，多半是把几个正交维度——加工方式、
控制方式、体量——全部塞进了 is-a。更稳的做法是：主干层次只保留
一个最本质的维度，其余维度用属性表达（\texttt{hasControlType}、
\texttt{hasSizeClass}），真需要「数控且大型」的集合时，
交给查询条件或定义类去拼。维度分解的功夫此刻多花一分，
3.7 节体检时就少修十处。

\subsection{第五步：定义属性}

类回答「领域里有什么概念」，属性回答「概念之间怎么连、概念身上记什么」。
第 2 章的两族划分在这里落地。\textbf{对象属性}连接两个个体：
\texttt{canProcess}（\texttt{Equipment} \(\rightarrow\) \texttt{Material}，能力关系）、
\texttt{locatedIn}（\texttt{Equipment} \(\rightarrow\) \texttt{WorkCell}，空间关系）、
\texttt{precedes}（\texttt{Process} \(\rightarrow\) \texttt{Process}，工序先后，
具有传递性）、\texttt{operates}（\texttt{Worker} \(\rightarrow\) \texttt{Equipment}）。
\textbf{数据属性}把个体连到字面值：\texttt{hasPower}（\texttt{xsd:float}，单位千瓦）、
\texttt{hasSerialNumber}（\texttt{xsd:string}）、
\texttt{hasStatus}（取值限定在 Running、Idle、Maintenance 的枚举）。

三个工程要点。其一，\textbf{每条属性都要写明定义域与值域}——这不只是文档礼貌，
推理器将拿它们做类型推断与抓错（第 4 章演示）。其二，\textbf{命名用动词或
动词短语的小驼峰形式}（\texttt{canProcess}、\texttt{locatedIn}），
读起来像句子，反推 CQ 时也顺。其三，属性的逻辑特性——传递性、函数性、逆属性——
此刻先用注释记下，OWL 落地写法第 4 章给出。
还有一条预防针：发现自己想建「空闲设备」这种类时停一停——
「空闲」更像属性值而非类，这桩公案的彻底了断在 3.7 节。

属性自身也可以组成层次：\texttt{canTurn}（能车削）与
\texttt{canMill}（能铣削）都可以声明为 \texttt{canProcess} 的子属性，
于是「3 号车床能车削铝合金」自动蕴含「3 号车床能加工铝合金」——
CQ1 的查询无须知道车削与铣削之分。子属性是新手最常遗忘的表达力：
它让属性清单像类层次一样有粗有细，粗粒度供通用查询，细粒度供专业查询，
而推理器负责在两层之间自动搭桥。

\subsection{第六步：定义属性约束}

属性框架立起来后，给它上规矩——约束（facet）把口头的业务规则固化成
机器可检查的公理：

\begin{codebox}
\codeline{\textcolor{cmtgray}{\#\ ==============================}}
\codeline{\textcolor{cmtgray}{\#\ Step\ 6.\ 定义属性约束（Facets）}}
\codeline{\textcolor{cmtgray}{\#\ ==============================}}
\codeline{\textcolor{cmtgray}{\#\ 基数约束：每台设备恰有一个序列号}}
\codeline{\textcolor{cmtgray}{\#\ \ \ Equipment\ \(\sqsubseteq\)\ =1\ hasSerialNumber}}
\codeline{\textcolor{cmtgray}{\#\ 值域约束：功率非负}}
\codeline{\textcolor{cmtgray}{\#\ \ \ hasPower\ 的取值\ \(\ge\)\ 0}}
\codeline{\textcolor{cmtgray}{\#\ 全称约束：数控机床的控制器都是数控系统}}
\codeline{\textcolor{cmtgray}{\#\ \ \ CNCMachine\ \(\sqsubseteq\)\ \(\forall\)hasControl.CNCControl}}
\codeline{\textcolor{cmtgray}{\#\ 不相交：Equipment\ \(\sqcap\)\ Material\ \(\sqsubseteq\)\ \(\bot\)}}
\codeblank{}
\end{codebox}
\color{black}

\assetsource{ch03}{ontology-101-process}

四种最常用的约束各有分工：\textbf{基数约束}「每台设备恰有一个序列号」
把台账纪律写进模型；\textbf{值域约束}「功率非负」拦住脏数据；
\textbf{全称约束}「数控机床的控制器都是数控系统」表达配置规范；
\textbf{不相交公理}「设备与材料不相交」则是第 2 章强调过的防混淆地基——
没有它，推理器永远不会发现「把一批铝棒登记成了设备」这种错误。

但要带着第 2 章的清醒来用约束：OWL 的约束\textbf{不是数据库的 CHECK}。
在开放世界假设下，「没登记冷却系统」不会立刻报错——推理器只当那是「尚不知道」。
所以 CQ7 的超期检查走的是查询层的局部封闭路线
（\texttt{FILTER NOT EXISTS}），公理负责定义「应当如何」，
查询负责巡检「实际如何」。公理与查询的这种分工，是贯穿第 4 章的工程暗线。

约束的数量也要克制。每条公理都是给未来变更立下的军令状：
约束越多，模型越「硬」，领域演化时要连带改动的地方就越多。
经验上的取舍是：把\textbf{业务红线}——违反即事故的规则，
如「高功率必配冷却」——写成公理；把\textbf{当前惯例}——
眼下如此、明年未必的做法——留在查询或文档层。
分不清某条规则是红线还是惯例？去问当初给这条 CQ 背书的那位专家。
这正是背书机制的又一次回报：每条知识都有一位可以追问的主人。

\subsection{第七步：创建实例}

骨架完工，开始填血肉。示范实例 \texttt{Lathe\_003} 这样登记：
它是 \texttt{CNCLathe} 的实例；序列号「EQ-2023-0117」
（\texttt{hasSerialNumber}）；功率 15.5 千瓦（\texttt{hasPower}）；
位于工作单元 A（\texttt{locatedIn} 指向 \texttt{WorkCell\_A}）；
能加工 6061 铝合金（\texttt{canProcess} 指向相应的材料个体）。

两个工程提醒。其一，\textbf{手工只建少量示范实例}，目的有二：跑通 CQ 验收、
检验设计手感——填实例时哪里写得别扭，哪里就埋着层次或属性的设计问题，
实例是设计的第一块试金石。大批量实例绝不该手敲，而应从设备台账、MES
等业务系统经数据管道自动导入，那是第 7 章知识图谱构建的正题。
其二，\textbf{实例命名规范从第一天就定下}（如「类型前缀加编号」：
\texttt{Lathe\_003}、\texttt{WorkCell\_A}），否则导入管道上线之日，
就是命名风格打架之时。

\subsection{循环而非瀑布：迭代与一致性检查}

七个步骤的编号给人瀑布错觉，实际的轨迹永远是螺旋：
创建实例时发现缺属性，回到第五步；CQ 验收不通过，回到第四步调整层次；
细化过程中发现复用的上游本体不合身，回到第二步重新谈判。
每一轮迭代的出口设两道检查：\textbf{声明 profile 内的语义检查}——通过
Semantica 场景确认已支持的检查没有失败，以及 \textbf{CQ 回归}——对具名正例、
单因反例和上一发布版 oracle 做精确比较。未绑定的 DL 或 OntoClean 检查保持红灯。

节奏上的建议是\textbf{小步快走}：每轮认领一小撮 CQ，
两三周内就要见到“在声明场景中能回答问题”的 package，让专家同时看到结果、
反例与未通过项；不以一个查询返回数据替代整包验收。
这正是敏捷软件开发的节奏感——也再次印证本章的主题：
方法论不是束缚创造力的枷锁，而是让创造力不脱轨的护栏。

迭代还要回答一个「停机问题」：什么时候算够？七步法自身不作答，
答案在 3.2 节——当前迭代认领的 CQ 全部验收通过、
且没有新的高优先级 CQ 入池，即可冻结并发布一个版本。
「完美的本体」并不存在，存在的只有「回答得了当前问题清单的本体」。
让清单而不是完美主义来叫停，是 CQ 体系送给工程师的一剂心理解药——
故事一里那三个月的精雕细琢，缺的正是这一声哨。

\section{重型流程：METHONTOLOGY}

七步法适合身手敏捷的小团队。但换一个场景：你是供应商，
要在合同框架下为某集团交付一套设备调度本体，
对方要求里程碑评审、完整的过程文档、可审计的变更记录和多年期维护承诺——
这时「轻」就成了「裸奔」。为这类场景准备的重型方法论是
METHONTOLOGY（Fern\'andez-L\'opez 与 G\'omez-P\'erez 等人于 1997 年提出），
它是最接近软件工程规范的本体构建方法：阶段明确、产物可交付、过程可检查。

\subsection{生命周期：三类活动}

METHONTOLOGY 把项目中的活动分成三类，各司其职。
\textbf{管理活动}贯穿全程：进度控制与质量保证，对应项目经理的工作面。
\textbf{开发活动}是主线，依次为五个阶段：\textbf{规格说明}产出规格说明书，
写明本体名称、目的、形式化级别、范围、知识源与用户
（包内教学资产第 2 节给出了车间调度本体样例，一页纸即可容纳）；
\textbf{概念化}把领域知识整理成一组与实现语言无关的中间表格——这是全方法的灵魂，
下一小节专讲；\textbf{形式化}把表格翻译成描述逻辑等形式模型；
\textbf{实现}将其落成 OWL 文件；\textbf{维护}处理上线后的变更请求并保持文档同步。
\textbf{支撑活动}同样贯穿全程：知识获取、集成（即复用）、评估、文档化与版本管理。
支撑活动的强度随阶段起伏——知识获取在前期最密集，评估则越到后期越重，
直到以验收收官。

五个开发阶段中，最容易被外行低估的是维护。本体的「需求」是领域知识本身，
而领域在变：新设备类型入场、工艺规程改版、组织架构调整，
都会化作变更请求涌向本体。METHONTOLOGY 把维护列为正式阶段并要求文档同步，
正是在提醒投资方：本体不是一锤子买卖，预算里必须给维护留出位置——
这一点将在 3.8 节的治理机制中落实成具体动作。

\subsection{灵魂产物：五张概念化表格}

概念化阶段要求在写下任何一行 OWL 之前，先产出五张中间表格。
先看包内教学资产的节选，再逐张解读：

\begin{codebox}
\codeline{\textcolor{cmtgray}{\#\ ==============================}}
\codeline{\textcolor{cmtgray}{\#\ 3.\ 概念化产物（Conceptualization）}}
\codeline{\textcolor{cmtgray}{\#\ ==============================}}
\codeline{\textcolor{cmtgray}{\#\ METHONTOLOGY\ 要求产出一组中间表格，先于任何OWL代码：}}
\codeblank{}
\codeline{\textcolor{cmtgray}{\#\ (1)\ 术语表（Glossary\ of\ Terms，节选）}}
\codeline{\textcolor{cmtgray}{\#\ \ \ 术语\ \ \ \ \ \ \ \ \ \ 类型\ \ \ \ \ \ 描述}}
\codeline{\textcolor{cmtgray}{\#\ \ \ ----------\ \ \ \ ------\ \ \ \ --------------------------------}}
\codeline{\textcolor{cmtgray}{\#\ \ \ Equipment\ \ \ \ \ 概念\ \ \ \ \ \ 车间中可执行加工/测量任务的装置}}
\codeline{\textcolor{cmtgray}{\#\ \ \ CNCLathe\ \ \ \ \ \ 概念\ \ \ \ \ \ 数控车床}}
\codeline{\textcolor{cmtgray}{\#\ \ \ canProcess\ \ \ \ 二元关系\ \ 设备具备加工某材料的能力}}
\codeline{\textcolor{cmtgray}{\#\ \ \ hasPower\ \ \ \ \ \ 属性\ \ \ \ \ \ 设备额定功率，kW，非负实数}}
\codeline{\textcolor{cmtgray}{\#\ \ \ Lathe\_003\ \ \ \ \ 实例\ \ \ \ \ \ 3号数控车床}}
\codeblank{}
\codeline{\textcolor{cmtgray}{\#\ (2)\ 概念分类树（Concept\ Taxonomy）}}
\codeline{\textcolor{cmtgray}{\#\ \ \ Equipment}}
\codeline{\textcolor{cmtgray}{\#\ \ \ \ \ ├──\ ProcessingEquipment}}
\codeline{\textcolor{cmtgray}{\#\ \ \ \ \ │\ \ \ \ \ └──\ CNCMachine\ ──\ CNCLathe\ /\ CNCMillingMachine}}
\codeline{\textcolor{cmtgray}{\#\ \ \ \ \ └──\ MeasurementEquipment}}
\codeblank{}
\codeline{\textcolor{cmtgray}{\#\ (3)\ 二元关系表（Binary\ Relations\ Table）}}
\codeline{\textcolor{cmtgray}{\#\ \ \ 关系名\ \ \ \ \ \ 定义域\ \ \ \ \ \ \ 值域\ \ \ \ \ \ \ 逆关系\ \ \ \ \ \ \ \ 特性}}
\codeline{\textcolor{cmtgray}{\#\ \ \ --------\ \ \ \ ---------\ \ \ \ -------\ \ \ \ ----------\ \ \ \ ------}}
\codeline{\textcolor{cmtgray}{\#\ \ \ canProcess\ \ Equipment\ \ \ \ Material\ \ \ processedBy\ \ \ —}}
\codeline{\textcolor{cmtgray}{\#\ \ \ locatedIn\ \ \ Equipment\ \ \ \ WorkCell\ \ \ contains\ \ \ \ \ \ 函数性}}
\codeline{\textcolor{cmtgray}{\#\ \ \ precedes\ \ \ \ Process\ \ \ \ \ \ Process\ \ \ \ follows\ \ \ \ \ \ \ 传递性}}
\codeblank{}
\codeline{\textcolor{cmtgray}{\#\ (4)\ 概念字典（Concept\ Dictionary，节选）}}
\codeline{\textcolor{cmtgray}{\#\ \ \ 概念\ \ \ \ \ \ \ \ 实例标识\ \ \ \ \ \ \ \ 属性\ \ \ \ \ \ \ \ \ \ \ \ \ \ \ \ \ \ 关系}}
\codeline{\textcolor{cmtgray}{\#\ \ \ --------\ \ \ \ ------------\ \ \ \ ------------------\ \ \ \ ----------------}}
\codeline{\textcolor{cmtgray}{\#\ \ \ Equipment\ \ \ serialNumber\ \ \ \ hasPower,\ hasStatus\ \ \ canProcess,\ locatedIn}}
\codeblank{}
\codeline{\textcolor{cmtgray}{\#\ (5)\ 实例表（Instance\ Table，节选）}}
\codeline{\textcolor{cmtgray}{\#\ \ \ 实例\ \ \ \ \ \ \ \ 所属概念\ \ \ \ 属性赋值}}
\codeline{\textcolor{cmtgray}{\#\ \ \ ---------\ \ \ --------\ \ \ \ --------------------------------}}
\codeline{\textcolor{cmtgray}{\#\ \ \ Lathe\_003\ \ \ CNCLathe\ \ \ \ hasPower=15.5;\ hasStatus=Idle}}
\codeblank{}
\codeline{\textcolor{cmtgray}{\#\ 产物的价值：领域专家不读OWL，但能逐行评审这些表格——}}
\codeline{\textcolor{cmtgray}{\#\ 概念化表格是专家与本体工程师之间的"共同语言"}}
\codeblank{}
\end{codebox}

\assetsource{ch03}{methontology-lifecycle}

\textbf{术语表}（Glossary of Terms）是七步法第三步那张词汇清单的「转正」：
每个术语被赋予类型（概念、二元关系、属性、实例）并配上一句自然语言描述。
它是其余四张表的索引，也是评审会上的「字典」。

\textbf{概念分类树}（Concept Taxonomy）是类层次的表格化呈现，
与七步法第四步同源；差别在于身份——它是独立交付物，
专家不必打开任何工具就能审它。

\textbf{二元关系表}（Binary Relations Table）逐行登记每条关系的定义域、值域、
逆关系与逻辑特性。这张表的妙处在「逆关系」与「特性」两列：
它们逼着设计者把话说全——\texttt{locatedIn} 的逆是什么？\texttt{precedes}
传递吗？七步法里这些问题靠自觉，这里靠表格不留死角。

\textbf{概念字典}（Concept Dictionary）换一个视角，以概念为行，
汇总它的实例标识、属性与关系，回答「关于 \texttt{Equipment}
我们到底记录哪些东西」——相当于每个概念的「户口页」。

\textbf{实例表}（Instance Table）登记示范实例及其属性赋值，
它同时就是验收时的种子数据。

五张表格的真正价值一句话可以说尽：\textbf{领域专家不读 OWL，
但能逐行评审这些表格}。它们是专家与本体工程师之间的共同语言——
故事三里缺席的那座桥。更要紧的是时机：评审发生在 OWL 编码\textbf{之前}，
概念错误在纸面上就被改掉。软件工程的老定律「缺陷发现得越早，修复代价越小」，
在知识工程里原样成立。

实际评审中，五张表的出场顺序也有讲究：先过术语表，确认「词」
没有遗漏与歧义；再过概念分类树，确认「骨架」符合专家直觉；
然后是二元关系表与概念字典，确认「连接」与「户口」；实例表压轴，
用具体例子反查前四张表。任何一张表的评审都可能把问题打回上一张——
这个小循环正是概念化阶段内部的迭代，与七步法的螺旋一脉相承。

\subsection{与软件工程对照}

METHONTOLOGY 的骨架几乎可以和软件工程逐项对表，
这也解释了为什么有软件过程经验的团队上手它最快：

\begin{center}
\begin{tabularx}{\linewidth}{@{}llX@{}}
\toprule
\textbf{METHONTOLOGY} & \textbf{软件工程对应} & \textbf{关键产物} \\
\midrule
规格说明 & 需求分析 & 规格说明书（目的、范围、知识源、用户） \\
概念化 & 概要设计 & 五张概念化表格 \\
形式化 & 详细设计 & 描述逻辑模型 \\
实现 & 编码 & OWL 文件 \\
评估 & 测试 & 一致性报告与 CQ 验收记录 \\
文档化 & 交付文档 & 标注齐全的本体加变更日志 \\
维护 & 运维 & 版本序列与变更记录 \\
\bottomrule
\end{tabularx}
\end{center}

评估一栏值得多看一眼：教学资产把评估拆成\textbf{一致性}、\textbf{能力}与
\textbf{质量}三个层面。当前 ch03 只原生绑定 CQ1 能力场景；完整 DL 一致性与
OntoClean 检查器并未实现。方法框架可以抄进项目验收单，绿色状态只能由实际证据填写。

对照表也暴露了 METHONTOLOGY 的年代局限：它诞生于瀑布模型盛行的年代，
原始表述里阶段感很重。当代实践早已把它「敏捷化」——
五个阶段在每轮迭代里小规模地全部走一遍，而不是一个阶段闷头三个月。
换句话说，METHONTOLOGY 的当代价值不在流程编排，而在\textbf{产物清单}：
无论节奏多快，规格说明书、五张表格、评估记录这些交付物本身，
依然是企业级项目的金标准。

\subsection{何时选择重型流程}

判断信号很直白：\textbf{多组织协作}（概念分歧需要书面裁决）、
\textbf{合同交付与审计要求}（过程必须留痕）、\textbf{生命周期以年计}
（人员会流动，知识必须落在文档而非脑袋里）——命中任意两条，就值得上重型。
反之，单团队内部的快速迭代项目用七步法即可，强行全套 METHONTOLOGY
只会生产没人读的文档。两者也并不互斥：常见的折中是以七步法为日常节奏，
借用 METHONTOLOGY 的产物模板补文档短板。若资源只够维护两张表，
多数团队会留下术语表与二元关系表——理由请你在思考题 5 里自己论证。

\section{场景化组合：NeOn 方法}

第三个现实是：今天几乎没有本体从真空中长出来。企业手里早就有分类标准、
数据字典、Excel 台账，公网上还有大量现成的本体与词表。
NeOn 方法（2009 年，源自同名欧洲研究项目）正视这一点，
干脆放弃「唯一正确流程」的想法，改为提供\textbf{九种可组合的场景}——
按你手头有什么资源，组装属于你的流程：

\begin{center}
\begin{tabularx}{\linewidth}{@{}lp{0.30\linewidth}X@{}}
\toprule
\textbf{场景} & \textbf{名称} & \textbf{做什么} \\
\midrule
场景 1 & 从头构建 & 无可复用资源，走规格说明、概念化到实现的主线 \\
场景 2 & 复用非本体资源 & 把术语表、分类标准、数据库模式等改造成本体 \\
场景 3 & 复用本体资源 & 引入或裁剪已有本体为己所用 \\
场景 4 & 复用并重构本体资源 & 对引入的本体改写调整，使之适配自身需求 \\
场景 5 & 复用并合并本体资源 & 同主题多个本体取长补短，合为一体 \\
场景 6 & 复用、合并并重构 & 合并之后再做深度改造 \\
场景 7 & 复用本体设计模式 & 套用「部分—整体」「角色」等成熟模式 \\
场景 8 & 重组本体资源 & 对既有本体做模块化、剪枝与扩展 \\
场景 9 & 本体本地化 & 把本体适配到新的语言与术语环境 \\
\bottomrule
\end{tabularx}
\end{center}

九个场景共享场景 1 的主干——需求规格照样要写，CQ 照样要过——
区别只在中段插入哪些资源加工活动。换句话说，NeOn 不是第十种流程，
而是一张「流程拼装说明书」。

工程现实中出场率最高的组合是\textbf{场景 2 加场景 3}。设想企业已有一份
《设备分类与编码标准》：三级类目、每级有编码与名称——典型的非本体资源。
场景 2 的加工分四步走。第一步\textbf{逆向工程}：把编码表读成一棵树，
一级类目「金属切削机床」下挂「车床」「铣床」，再往下是「数控车床」「仪表车床」。
第二步\textbf{语义审查}，也是最考验功力的一步：逐条追问「下级真的是上级的
\textbf{一种}吗」。分类标准是给管理用的，常把「按用途分」「按结构分」混在一棵树里，
还可能混进采购口径（「进口设备」「租赁设备」）——这些都不是真正的 is-a 关系，
照搬进类层次就埋下了 3.7 节要清算的隐患。第三步\textbf{转换}：
通过审查的层级翻译成 \texttt{rdfs:subClassOf}，类目名称进 \texttt{rdfs:label}，
原编码保留为数据属性（如 \texttt{hasLegacyCode}，留作与老系统对账的脐带），
可疑的层级关系降级为普通属性。第四步\textbf{补全}：对照 CQ 清单，
补上分类标准里根本没有的属性与公理——分类标准只关心「是什么」，
不关心「能加工什么」。

这段加工里最值得复盘的是第二步，给一个典型的反面教材：
某分类标准把「数控设备」与「车床」并列为二级类目，
于是「数控车床」在两处各出现一次、编码还不相同，
照搬的结果是一台设备两个身份，台账永远对不上。
本体化时的正确处理是把「数控」识别为控制方式维度，
用属性 \texttt{hasControlType} 表达，类树上只保留「车床」一支——
这正是 3.3.4 小节「维度分解」原则与 3.7 节体检思想的提前上演。

然后场景 3 登场：把转换得到的设备类挂接到 IOF 对应概念之下，
计量单位对齐 QUDT。完成这一步，企业的私有分类就拿到了与外部世界互操作的接口。
注意终点没有变：无论场景怎么组合，产出的本体仍要走 3.2 节的 CQ 验收——
\textbf{场景改变的是路径，不是验收门}。

九个场景里还有一位值得单独引荐的配角：场景 7「复用本体设计模式」。
设计模式（Ontology Design Pattern，ODP）是比整本本体更小的复用单元——
「部分—整体」「角色扮演」「参与者—过程」这些反复出现的建模难题，
社区早已沉淀出经过验证的标准解法。3.7 节将要展开的「类型与角色分离」
就是其中最经典的一例。把模式当乐高积木，比整本搬运别人的本体更轻、
更可控，特别适合在七步法的第四、五步随用随取。

\section{站在肩膀上：本体复用工程}

NeOn 场景 3 的一句「复用本体资源」，背后是一整套工程动作。
复用做对了是杠杆——别人十年打磨的概念体系一夜为你所用；
做错了是负债——语义错配与上游变更会折磨项目的整个余生。
本节把复用拆成\textbf{找、评、用}三步，并把陷阱标在地图上。
全部资源的访问入口集中在附录 A，本节只讲方法。

\subsection{找：去哪里搜本体}

入口分三类。\textbf{综合目录}：LOV（Linked Open Vocabularies）收录通用词表，
支持按词条全文检索，适合「先看看有没有人定义过 \texttt{Machine}」式的初筛。
\textbf{领域本体库}：BioPortal 是生物医学领域规模最大的本体门户之一；
OBO Foundry 以严格的共建原则著称——一个领域一个权威本体、共享顶层、开放许可，
即便你不做生物医学，它的治理模式也值得借鉴。\textbf{标准组织与基金会}：
ISO/IEC 21838 系列（顶层本体框架与 BFO）、工业本体基础 IOF、
W3C 推荐标准（Time、SOSA 等）。检索技巧只有一条但很关键：
\textbf{用 CQ 里的核心名词（英文）去搜，命中后先读定义与公理，再看类名}——
类名相同不等于概念相同。

检索也可能空手而归，而这同样是有价值的结论：它意味着你要建的
可能确实是空白地带，自建的投入因此有了正当性。记得把
「查过哪里、候选有谁、为什么不采用」写进规格说明书——
三年后的接手者会感谢这段记录，因为他不必再把整个互联网重新搜一遍。

\subsection{评：四个维度打分}

候选本体到手，逐一过四道关：

\begin{center}
\begin{tabularx}{\linewidth}{@{}lXX@{}}
\toprule
\textbf{维度} & \textbf{核心问题} & \textbf{检查手段} \\
\midrule
覆盖度 & 我的术语表有多少能在其中找到对应？ & 拿第三步术语表逐条对照，统计命中与缺口 \\
一致性 & 它的顶层承诺、建模风格与我相容吗？ & 查它对齐的顶层本体；推理器跑一遍看有无不可满足类 \\
活跃度 & 它还活着吗？ & 看最近发布日期、变更日志、问题响应速度与维护组织 \\
许可 & 我能在商业产品中使用、修改、再分发吗？ & 逐字读许可证条款，确认义务与限制 \\
\bottomrule
\end{tabularx}
\end{center}

一致性维度最容易被低估：两个分别对齐不同顶层本体的资源，
其「过程」「对象」的切分方式可能根本无法调和，硬拼的代价远高于自建。
活跃度则更多关乎风险而非质量：一个多年没有更新的本体，内容也许仍然正确，
但它意味着没有人替你修缺陷、回议题——引入它，约等于把维护责任一并买断。
评估结论未必是「用或不用」的二选一——更常见的结论是「用它的某个模块」
或「只对齐它的骨干概念」，这就引出第三步的姿势选择。

\subsection{用：三种引入姿势}

\textbf{姿势一：\texttt{owl:imports} 全量引入。}一行导入声明，
对方整个本体连同全部公理进入你的逻辑空间。优点是省事且随上游演进；
代价是推理负担全盘继承，且上游一旦大改，你被动跟随。
适合稳定、精悍的小型本体（如 W3C Time）。

\textbf{姿势二：概念映射。}本地自建类，再用 \texttt{rdfs:subClassOf} 或
\texttt{owl:equivalentClass} 与外部概念对齐。耦合最松，自家本体保持苗条，
推理范围可控；代价是对齐断言的正确性全靠自己负责，上游更新也不会自动传导。
适合只需「挂顶层」的场景——3.5 节对齐 IOF 用的正是这一式。

\textbf{姿势三：MIREOT 式摘取。}介于两者之间：只把需要的术语
（IRI、标签、定义及最小上下文）抽进自家本体，不背全部公理。
MIREOT（Minimum Information to Reference an External Ontology Term）
是 OBO 社区的成熟惯例，有现成工具支持。它保留了外部 IRI——
将来与他人交换数据时，「同一个概念用同一个标识符」的互操作红利仍然在手。

三种姿势怎么选？一条实用的决策线：上游小而稳、整体都要——全量引入；
只需借骨架对齐、保持自家苗条——概念映射；想要的术语散落在几个大部头里——
MIREOT 摘取。同一个项目对不同上游混用三种姿势，不但正常，
而且往往是最优解：Time 全量引入、IOF 概念映射、某个庞大领域本体里
只摘三个类，各得其所。

\subsection{复用的陷阱}

\textbf{陷阱一：语义不合。}同名不同义是常态——「Process」在业务流程建模语境
与在 BFO 语境下几乎是两个世界；「Resource」在调度语境与在 Web 架构语境
也是貌合神离。只看类名就引入，等于没读条款就签合同。解药是笨功夫：
读定义、读公理、看它的实例长什么样。

\textbf{陷阱二：过度依赖。}上游一次重构，下游集体断链。
解药是把上游当成「受管依赖」：锁定版本 IRI 引入而不是追最新版，
把升级当成显式决策；上游变更后，先在自家 CQ 回归上跑一遍再合入。

\textbf{陷阱三：许可与治理遗忘。}引入即承担其许可义务，
商用前务必确认修改权与再分发条款；引入的概念也要纳入自家变更治理，
不能「进了门就没人管」。

最后一条原则给所有姿势收尾：\textbf{复用不豁免验收}。
引来的类同样要过 CQ 回归、同样要接受 OntoClean 体检——
「权威出品」只说明它适合作者的场景，不自动等于适合你的场景。

把三步法压缩成一张随身检查单：\textbf{找}——列出至少三个候选并记录出处；
\textbf{评}——覆盖度有对照数字、一致性有推理报告、活跃度有日期、许可有结论；
\textbf{用}——写明引入姿势与版本号，对齐断言单独成文件以便回滚。
最后为每个上游建立一条「升级例行程序」：订阅其发布通知，
升级前在分支上跑全量 CQ 回归。做到这一步，
复用就不再是一次性的拿来主义，而是一段被妥善管理的长期关系。

\section{给类层次体检：OntoClean}

类层次建好了，推理器一片绿灯，可评审会上总有人嘀咕：
「设备挂在资源下面，总觉得哪里怪怪的。」这种「怪怪的」直觉非常宝贵——
推理器检查的是\textbf{逻辑一致性}（公理之间有没有矛盾），
而「设备是一种资源」在逻辑上毫无矛盾；它的毛病出在\textbf{概念合理性}上：
把「类型」与「角色」搅在了一起。OntoClean（Guarino 与 Welty，2002 年）
正是为这类问题而生的体检方法，思路独辟蹊径：\textbf{不看类的内容，
只看类的形式性质}——所谓元属性（metaproperty）——再用少量规则机械地揪出混淆。
它是对第 2 章类层次知识的一次深水区延伸。

\subsection{四个元属性}

\textbf{刚性}（Rigidity）问的是：实例会不会在不失去自身的前提下脱离这个类？
\textbf{+R}（刚性）：实例永远是该类实例——一台车床在报废拆解之前永远是设备，
它不可能「暂时不当设备」，所以 \texttt{Equipment} 是刚性类。
\textbf{\(\sim\)R}（反刚性）：所有实例都\textbf{可能}不再属于该类——
每台「空闲设备」都会重新忙起来，「空闲」描述的是状态而非本质，
\texttt{IdleEquipment} 是反刚性类。还有居中的 \textbf{-R}（非刚性）：
部分实例可变——「坚硬的东西」里有些会软化、有些不会。
判定口诀一句话：\textbf{会「变回去」的，不是刚性类}。

\textbf{同一性}（Identity）问的是：这个类有没有携带「判定是不是同一个」的判据？
\texttt{Equipment} 是 \textbf{+I}：序列号就是判据——序列号相同即同一台设备。
「红色的东西」是 \textbf{-I}：两团红色之间不存在统一的「同一红物」判定标准，
红色这个类不负责回答个体同一性问题。同一性判据往往是「真类型」的指纹：
能说清楚怎么数、怎么认的概念，多半是层次里的正经骨干。

\textbf{统一性}（Unity）问的是：该类的实例有没有「一个整体」的边界判据？
\texttt{Machine} 是 \textbf{+U}：一台机器边界分明，哪些零件属于它说得清楚。
\texttt{Metal} 是 \textbf{-U}：一摊金属没有「一个金属」的天然边界——
物质名词普遍如此。统一性失配的高发地带，正是「对象」与「构成它的材料」
被混为一谈的地方。

\textbf{依赖性}（Dependence）问的是：该概念的成立是否依赖其他实体的存在？
\texttt{Operator}（操作员）是 \textbf{+D}：没有被操作的设备，谈不上操作员；
\texttt{SparePart}（备件）也是 +D：备件之为备件，是相对某台整机而言的。
依赖性是识别「角色」的雷达——角色总是依赖某个语境或相对方才成立。

四个属性合起来给出一组朴素的直觉：\textbf{「类型」（+R，常带 +I+U）回答
「它是什么」}——车床永远是车床；\textbf{「角色与状态」（\(\sim\)R，常带 +D）回答
「它当前充当什么」}——空闲、瓶颈、备件。类型适合做类层次的骨架，
角色与状态适合用属性表达。这一句，几乎就是 OntoClean 的全部精神。

标注时的两个常见疑难也提前拆解。疑难一：「一个类能既 +R 又 +D 吗？」
能——四个元属性彼此独立，逐个判断即可；真正要避免的是对着类名空想，
正确姿势是想着它的\textbf{典型实例}提问：拿 3 号车床来问刚性，
拿一批 6061 铝棒来问统一性。疑难二：「\texttt{Material} 到底 +R 还是
\(\sim\)R？一块铝锭被铸成零件之后还算不算材料？」这类边界案例
往往没有唯一正确答案，而裁决本身就有文档价值——把裁决理由写进
\texttt{rdfs:comment}，比裁决结论更重要。元属性标注的过程价值
常常大于结果：它强迫团队把「我们说的设备到底指什么」这种
从未明说的共识，第一次摆上桌面。

\subsection{标注流程与标注表}

工程上的标注流程五步走：(1) \textbf{选骨干类}——不必全量标注，
挑层次中承重的二三十个类即可；(2) \textbf{逐类四问}——会变回去吗？
有判同标准吗？有「一个」的边界吗？依赖他者吗？(3) \textbf{填写标注表}；
(4) \textbf{扫描父子边}——用下一小节的三条规则检查每一条 \(\sqsubseteq\)；
(5) \textbf{修正复扫}，直至无违规。制造本体骨干类的标注结果如下：

\begin{codebox}
\codeline{\textcolor{cmtgray}{\#\ ==============================}}
\codeline{\textcolor{cmtgray}{\#\ 2.\ 制造本体标注示例}}
\codeline{\textcolor{cmtgray}{\#\ ==============================}}
\codeline{\textcolor{cmtgray}{\#\ \ \ 类\ \ \ \ \ \ \ \ \ \ \ \ \ \ \ \ \ \ \ \ \ \ R\ \ \ \ \ I\ \ \ \ \ U\ \ \ \ \ D\ \ \ \ 说明}}
\codeline{\textcolor{cmtgray}{\#\ \ \ --------------------\ \ \ \ ---\ \ \ ---\ \ \ ---\ \ \ --\ \ \ ------------------}}
\codeline{\textcolor{cmtgray}{\#\ \ \ Equipment\ \ \ \ \ \ \ \ \ \ \ \ \ \ \ +R\ \ \ \ +I\ \ \ \ +U\ \ \ \ -D\ \ \ 本质类型}}
\codeline{\textcolor{cmtgray}{\#\ \ \ CNCLathe\ \ \ \ \ \ \ \ \ \ \ \ \ \ \ \ +R\ \ \ \ +I\ \ \ \ +U\ \ \ \ -D\ \ \ 本质类型}}
\codeline{\textcolor{cmtgray}{\#\ \ \ Material\ \ \ \ \ \ \ \ \ \ \ \ \ \ \ \ +R\ \ \ \ +I\ \ \ \ -U\ \ \ \ -D\ \ \ 物质，无统一边界}}
\codeline{\textcolor{cmtgray}{\#\ \ \ IdleEquipment\ \ \ \ \ \ \ \ \ \ \ \textasciitilde{}R\ \ \ \ +I\ \ \ \ +U\ \ \ \ -D\ \ \ 状态（临时）}}
\codeline{\textcolor{cmtgray}{\#\ \ \ BottleneckResource\ \ \ \ \ \ \textasciitilde{}R\ \ \ \ +I\ \ \ \ +U\ \ \ \ +D\ \ \ 角色（依赖调度上下文）}}
\codeline{\textcolor{cmtgray}{\#\ \ \ Operator\ \ \ \ \ \ \ \ \ \ \ \ \ \ \ \ \textasciitilde{}R\ \ \ \ +I\ \ \ \ +U\ \ \ \ +D\ \ \ 角色（人担任）}}
\codeline{\textcolor{cmtgray}{\#\ \ \ SparePart\ \ \ \ \ \ \ \ \ \ \ \ \ \ \ \textasciitilde{}R\ \ \ \ +I\ \ \ \ +U\ \ \ \ +D\ \ \ 角色（相对整机而言）}}
\codeblank{}
\codeline{\textcolor{cmtgray}{\#\ 直觉：}}
\codeline{\textcolor{cmtgray}{\#\ \ \ "类型"（+R）回答"它是什么"——车床永远是车床}}
\codeline{\textcolor{cmtgray}{\#\ \ \ "角色/状态"（\textasciitilde{}R）回答"它当前充当什么"——空闲、瓶颈、备件}}
\codeblank{}
\end{codebox}
\color{black}

\assetsource{ch03}{ontoclean-evaluation}

这张表值得逐行咀嚼。前两行 \texttt{Equipment}、\texttt{CNCLathe}
是教科书式的本质类型（+R +I +U）：永远是它、认得出它、数得清它。
\texttt{Material} 的 \textbf{-U} 是警示灯：材料是物质而非对象，
它与设备的关系只能是「构成」而非「继承」。后四行清一色 \(\sim\)R——
「空闲」是状态，「瓶颈」「操作员」「备件」是角色；
其中 \texttt{BottleneckResource} 的 +D 说得很清楚：
「瓶颈」只有在特定排程语境下才成立，换一张排程表它就不是瓶颈了。
一行行读下来，层次里\textbf{谁是骨、谁是皮}，一目了然。

这张表还有一个隐藏用法：当作新成员的入职试卷。请新加入的工程师
独立标注同一批类，再与团队既有标注对差异——分歧点要么说明新人
还没吃透领域，要么说明团队文档没把判据写清楚，无论哪种都值得当场补课。
比起通读三天文档，一张三十行的标注表更能让人理解这个本体「为什么长这样」。

\subsection{三条规则与三类典型违规}

标注完成后，检查只剩机械劳动。三条核心规则：
\textbf{规则一（刚性约束）：刚性类不能以反刚性类为父类}——
+R \(\sqsubseteq\) \(\sim\)R 即违规。直觉论证：若 \texttt{Equipment}
挂在反刚性的 \texttt{Resource} 之下，那么「每台设备都可能不再是资源」，
顺着继承就得承认「每台设备都可能暂时不当设备」——荒谬。
\textbf{规则二（同一性约束）：子类必须继承父类的同一性判据，
携带不相容判据的类之间不能有包含关系}——「按序列号判同」的类
不能挂在「按重量批次判同」的类下面。
\textbf{规则三（统一性约束）：有统一性的类不能以反统一性类为父类}——
「一台」不能是「一摊」的子类。

三条规则对应三类高发违规。\textbf{违规一：角色当类型。}
最常见也最隐蔽，「设备是资源」听起来天经地义：

\begin{codebox}
\codeline{\textcolor{cmtgray}{\#\ 违规示例1：}}
\codeline{\textcolor{cmtgray}{\#\ \ \ Equipment(+R)\ \(\sqsubseteq\)\ Resource(\textasciitilde{}R)}}
\codeline{\textcolor{cmtgray}{\#\ \ \ "设备是资源"——但\ Resource\ 是调度语境下的角色（反刚性），}}
\codeline{\textcolor{cmtgray}{\#\ \ \ 设备是刚性类型，刚性类挂在反刚性类下，违反规则1}}
\codeblank{}
\codeline{\textcolor{cmtgray}{\#\ 修正：用"角色扮演"关系替代继承}}
\codeline{\textcolor{cmtgray}{\#\ \ \ Equipment\ playsRole\ ResourceRole}}
\codeline{\textcolor{cmtgray}{\#\ \ \ （类型层次与角色层次分离，角色通过对象属性关联）}}
\end{codebox}

\assetsource{ch03}{ontoclean-evaluation}

修正后的图景值得记住：\textbf{类型层次与角色层次分离}，两者用对象属性
\texttt{playsRole} 桥接。调度模块面向 \texttt{ResourceRole} 编程，
设备本体保持稳定——角色再多变，骨架不动摇。

\textbf{违规二：状态当子类。}把 \texttt{IdleEquipment} 之类的状态类
塞进继承层次，意味着车床每次开动都要「退出」该类、停机又「加入」——
而 OWL 的类成员资格不是为高频进出设计的，实例迁移既昂贵又易错。
修正：状态用数据属性表达——\texttt{Lathe\_003 hasStatus Idle}，状态变化只改一个属性值。
回看 3.2 节的 CQ2「当前哪些设备空闲」：它从头到尾不需要什么
\texttt{IdleEquipment} 类，一条按 \texttt{hasStatus} 过滤的查询足矣。
CQ 驱动设计与 OntoClean 在这里殊途同归——前者从需求侧、后者从质量侧，
共同否决了同一个过度建模。顺带提防它的变体「枚举状态当子类」：
有团队把每种状态都建成类（空闲设备、运行设备、维护中设备），
层次瞬间膨胀三倍，推理器还要为每次状态切换支付实例重分类的代价。
状态机请留在属性与数据层——这是本节第二次重复这句话，
因为它是实践中复发率最高的病。

\textbf{违规三：构成当继承。}「设备是金属」混淆了「\textbf{由}金属构成」
与「\textbf{是}金属」：\texttt{Equipment} 是 +U，\texttt{Metal} 是 -U，
正撞规则三的枪口。修正：用构成关系 \texttt{madeOf} 连接。
一个快速嗅探法：凡是改用「由……做成」「含有」复述反而更自然的「is-a」，
都值得立案审查。

\subsection{何时体检：局部审查策略}

对全本体做四元属性完整标注，成本高且收益递减——OntoClean 在工程中
的正确打开方式是\textbf{局部审查}，推荐三个时机：
(1) \textbf{类层次初版评审时}：只标注刚性一项，就能揪出「角色当类型」
这类发生率最高的错误，投入产出比最佳；
(2) \textbf{两个本体合并时}：重点核对双方的同一性判据是否冲突——
规则二的主战场，合并前不查，合并后必痛；
(3) \textbf{推理或查询出现反直觉结果时}：顺着结果反查涉事类的元属性，
往往能把「概念混淆」从一堆公理里精确定位出来。

包内教学资产最后给出一条经验法则，值得带进评审会验证：
\textbf{会「变回去」的概念——状态、角色、阶段——不要放进
\texttt{rdfs:subClassOf} 层次}，用 \texttt{hasStatus}、\texttt{playsRole}、
\texttt{hasPhase} 这样的属性表达。第 2 章打好的类层次地基，
配上这条法则与三条规则，足以让你的 TBox 在多数评审中昂首过关。

最后把 OntoClean 放回全章版图：CQ 管「建得有没有用」，
推理器管「建得有没有矛盾」，OntoClean 管「建得讲不讲理」。
三者查的是三种不同的病，谁也替代不了谁——这正是 3.4 节
评估三层面（能力、一致性、质量）的另一种说法。
一个只跑推理器的团队，等于体检只量了血压。

\section{一群人怎么建：团队协作与治理}

前几节解决的是「一个人怎么把本体建对」；真实项目里更难的问题是
「一群人怎么长期共建而不打架」。本体是共享的概念契约，
契约的维护本身就需要制度。本节给出四件治理装备：角色分工、版本管理、
评审会与文档底线。它们不如前几节「学术」，却往往决定项目的生死。

\subsection{三种角色}

成熟的本体团队至少要分清三顶帽子：

\begin{itemize}
  \item \textbf{本体工程师}：负责建模决策、公理编写、工具链与质量体检；
        交付物是本体文件、概念化表格与变更日志。是方法论的执行者。
  \item \textbf{领域专家}：负责提供知识、评审概念化表格、为每条 CQ
        给出标准答案。不要求会 OWL——这正是 3.4 节坚持表格评审的原因；
        专家的时间是项目最稀缺资源，要花在评审与判例上，而不是学语法上。
  \item \textbf{数据工程师}：负责实例层——台账与 MES 数据的映射、清洗、
        导入管道（第 7 章）与实例质量监控。本体上线后的日常健康大半系于此。
\end{itemize}

三顶帽子之外，还需要一位有决策权的负责人为范围分歧拍板。
小团队当然可以一人兼多角，但有一条红线：\textbf{「领域专家评审」与
「工程师自检」两顶帽子不能戴在同一个头上}——否则就回到了故事三：
自己出题自己改卷，闭门造车而不自知。

角色清单之外，还有一张「时间表」值得在开工前就谈定：
领域专家每轮迭代投入多少小时（评审一次、答疑一次，通常两到四小时足矣）、
答疑请求多久内响应、标准答案的最终拍板权在谁。
把专家时间白纸黑字写进项目计划，远比口头的「需要时随叫随到」可靠——
故事三的专家缺位，一半是工程师没去请，
另一半往往是专家的档期从未被正式预留。

\subsection{版本管理：语义版本号与变更日志}

本体和代码一样会演化，但版本号只是索引，不能代替内容身份与回归。
把“兼容性”翻译到知识层面：\textbf{删除类、改变概念语义}等
不兼容变更升主版本——下游的查询与映射可能因此失效，必须显式通告；
\textbf{新增类与属性}升次版本——旧查询照跑不误；
\textbf{修订标签、注释}通常升修订号；但最终身份还要绑定 package manifest、
asset hash、runtime commit/artifact 与 snapshot/diff。每个版本的\textbf{变更日志}
至少回答：改了什么、为什么改、影响哪些 CQ、由谁裁决、哪份 receipt 可以复核。

把 3.2 节与 3.7 节的成果接进来，就得到一道完整的\textbf{发布门禁}：
声明 profile 内的检查、CQ 正反例回归、需要的 OntoClean/SHACL 检查、文档与来源证据
全部完成才允许发布。任何 absent/partial/unsupported 都必须保留为阻断理由。
本体、CQ、规则、夹具、场景与变更记录在同一个 Semantica package 中共同版本化，
任何 receipt 都能回答“改了什么、用什么运行、有没有破坏旧能力”。

分支策略可以照搬代码仓库的玩法：可发布基线必须门禁全绿；开发中的 package
允许显式 blocked，但不得冒充 release。大重构在分支上进行，合并前跑全量回归。
还要预警一个不同点：
多人同时编辑同一个本体文件的冲突远比代码难合——语法层面的三方合并
对 OWL 文件几乎不可用。更现实的协作方式是\textbf{按模块分工}：
类层次、属性、实例分属不同文件，用 \texttt{owl:imports} 组装成整体，
谁的模块谁做主，跨模块改动走评审。模块边界本身，
也顺势成了团队分工的边界。

\subsection{评审会怎么开}

评审会是专家参与的主战场，开法有讲究。\textbf{节奏}：每轮迭代末开一次，
控制在一到两小时——拖到三小时，后半场的「同意」都是疲劳驾驶。
\textbf{材料}：概念化表格的增量部分、变更日志、CQ 通过情况看板；
铁律一条——\textbf{不要投影 OWL 源码}，那等于邀请专家集体沉默。
\textbf{流程}：逐行过表格，专家用业务语言挑战（「我们不这么叫」
「这两个其实是一回事」「这种情况漏了」），工程师当场记录而不当场辩论；
新冒出来的问题改写成候选 CQ 入池；一时争不清的标成待办，
留给负责人会后裁决。\textbf{产出}：问题清单、决议记录与新增 CQ。
再送一个屡试不爽的小技巧：评审会用「跑给你看」开场——
现场执行三条上轮承诺的 CQ 查询，让专家亲眼看到自己上次提的问题
已经能被回答。信任就是这样一轮一轮积累起来的。
线下会议之外还可以加一条异步通道：把概念化表格放进共享文档，
专家平时随手批注，评审会只集中处理批注中的争议项——
异步消化掉七八成琐碎问题，宝贵的会议时间才轮得到真正的概念分歧。

\subsection{文档的最低标准}

不设上限，但要守底线。\textbf{底线一}：每个类与每条属性都必须有
中英双语的 \texttt{rdfs:label}，以及一句自然语言定义形式的
\texttt{rdfs:comment}——定义最好写出判据：什么算、什么不算。
反面教材是只有 IRI 的「裸类」：三个月后连作者自己都想不起
\texttt{TempEquipment} 当初指什么。双语标签也不只是国际化礼貌：
中文标签供专家评审，英文标签供与国际本体对齐，
它们是 3.4 节「共同语言」原则在文档层的延续。
\textbf{底线二}：本体库根目录放一页 README，写明范围声明、命名规范、
复用清单（引了谁的哪个版本）与版本对照表。
\textbf{底线三}：变更日志不断档。三条底线加起来不过几人日的投入，
却是三年后接手者眼里项目「还活着」与「已成遗迹」的分水岭。

工具可以让守底线的成本趋近于零：把「无标签类检测」「无注释属性检测」
写成脚本挂进发布门禁，裸类根本提交不进主干；变更日志由提交说明自动汇总；
README 里的版本对照表随发版脚本自动更新。治理的最高境界不是指望人人自觉，
而是让不合规的东西无路可走——这句话同样适用于 CQ 回归与 OntoClean 抽查。

\section{本章小结}

把全章收拢成一张行军图。\textbf{一条主线}：能力问题贯穿始终——
它划范围（治过度建模）、驱动设计（反推类与属性）、定验收（治「永远快好了」），
做成自动回归后，知识工程第一次有了软件工程式的安全网。
\textbf{两套流程}：轻量的 Ontology 101 七步法以循环迭代见长，
适合中小团队快速起步；重型的 METHONTOLOGY 以五张概念化表格为灵魂，
让不读 OWL 的专家全程在场，适合企业级交付。NeOn 则把流程拆成九个可组合场景，
回应「从已有资源出发」的工程常态。\textbf{一个杠杆}：本体复用按「找、评、用」
三步推进，三种引入姿势各有代价，且复用不豁免验收。
\textbf{一次体检}：OntoClean 用四个元属性与三条规则给类层次做质检，
精髓一句话——类型进层次，角色与状态进属性。
\textbf{一套治理}：角色分工、语义版本、评审会与文档底线，
让方法论在多人长期协作中持续生效。

这五件装备的次序并非随意：范围错了，流程再严也只是高效地走错路；
流程对了，质量体检才有对象；质量稳了，治理机制才有值得守护的资产。
所以排查问题时永远从上游开始——先怀疑 CQ，再怀疑层次，最后才怀疑文档。

方法论给出了路线图，下一步需要交通工具。第 4 章将把本章反复预告的
OWL、Protégé 与 SPARQL 正式请上场——你在 3.2 节写下的验收查询，
只有在 Semantica 中绑定输入与 oracle 后才成为可复算场景；其余仍是待实现规格。

如果只许用一句话向同事转述本章，可以这样说：
\textbf{先写问题再建模型，先有表格再有代码，先做体检再上线，
先定规矩再多人开工}。四个「先」字，恰好把三个失败故事各自反过来讲了一遍。

\subsection*{思考题}

\begin{enumerate}
  \item 回到 3.1 节的三个失败故事，分别指出本章哪件工具能在最早的环节
        拦住它，并说明拦截机制各是什么。
  \item 把模糊需求「掌握车间设备的整体情况」改写成三条合格的能力问题，
        标注每条所属的类别（基础查询、关联推理、约束校验、统计分析），
        并为每条写出可操作的验收判据。
  \item 为 CQ8「哪些设备的维护记录超期」反推完整的本体元素清单
        （类、属性、约束），指出其中哪一环超出了纯 OWL 推理的能力范围，
        应当放在哪一层解决。
  \item 用 3.3.4 小节的三条自查规则审查如下类层次片段：
        \texttt{Equipment} 只有一个子类 \texttt{CNCMachines}，
        而 \texttt{CNCMachines} 之下又挂着子类 \texttt{Lathe\_003}。
        指出全部问题并给出修正方案。
  \item METHONTOLOGY 的五张概念化表格如果只许保留两张，你保留哪两张？
        请从「专家可评审性」与「信息覆盖面」两个角度论证取舍。
  \item 你所在企业已有一份三级设备分类编码标准，希望转成本体并对齐 IOF。
        写出你会采用的 NeOn 场景组合与步骤顺序，并指出哪一步最容易
        把「非 is-a」关系混进类层次、应如何识别。
  \item 用四个元属性标注下列类：\texttt{Person}、\texttt{Operator}、
        \texttt{Equipment}、\texttt{SparePart}、\texttt{Steel}。
        据此判断 \texttt{Operator \(\sqsubseteq\) Person} 与
        \texttt{SparePart \(\sqsubseteq\) Equipment} 是否违反三条检查规则，
        并说明各自应当如何建模才最稳妥。
  \item 为你的团队设计一份「本体版本发布门禁清单」：列出发布前必须通过的
        检查项、每项的执行方式与不通过时的处置，至少覆盖一致性、能力、
        质量与文档四个层面（提示：综合 3.2、3.7、3.8 三节）。
\end{enumerate}

\subsection*{本章包资产与场景指引}

下列 asset 均位于 \texttt{semantica.chapter\_packages.vol1.ch03}。

\begin{center}
\begin{tabularx}{\linewidth}{@{}>{\ttfamily\footnotesize}p{0.36\linewidth}X@{}}
\toprule
\normalfont\textbf{Semantica asset} & \textbf{内容与阅读建议} \\
\midrule
competency-questions & 十个 CQ 的生命周期：清单、分类、反推、验收查询与迭代；只有 CQ1 已绑定原生验收场景 \\
ontology-101-process & 七步法走查与常见错误；当前是工程规则材料，不是阶段门禁 runner \\
methontology-lifecycle & 生命周期、规格样例与概念化表格；当前是工程规则材料 \\
ontoclean-evaluation & 元属性与违规修正；自动 OntoClean 检查器仍缺失 \\
OE-V1-CH03-SCN-CQ-ACCEPTANCE-001 & CQ1 正例/单因反例的精确多重集验收；不等于 package release \\
\bottomrule
\end{tabularx}
\end{center}
